
\documentclass{article}
\usepackage{colortbl}
\usepackage{makecell}
\usepackage{multirow}
\usepackage{supertabular}

\begin{document}

\newcounter{utterance}

\centering \large Interaction Transcript for game `bbh', experiment `boolean\_expressions', episode 71 with reflexion{-}llama{-}3.1{-}8b{-}instruct.
\vspace{24pt}

{ \footnotesize  \setcounter{utterance}{1}
\setlength{\tabcolsep}{0pt}
\begin{supertabular}{c@{$\;$}|p{.15\linewidth}@{}p{.15\linewidth}p{.15\linewidth}p{.15\linewidth}p{.15\linewidth}p{.15\linewidth}}
    \# & \multicolumn{2}{c}{Player} && \multicolumn{2}{c}{Game Master} \\
    \hline

    \theutterance \stepcounter{utterance}  
    & & & \multicolumn{4}{p{0.6\linewidth}}{
        \cellcolor[rgb]{0.9,0.9,0.9}{
            \makecell[{{p{\linewidth}}}]{
                \texttt{\tiny{[P1$\langle$GM]}}
                \texttt{You have attempted this same task before and did not succeed. These are the plans you wrote after those attempts.} \\
\texttt{Trial \#1: New plan: } \\
\\ 
\texttt{In my previous attempt, I correctly identified the order of operations (PEMDAS) and applied it to the given Boolean expression. However, I failed to recognize that the expression "not not not not not not not True" can be simplified using the double negation law, which states that "not not P" is equivalent to "P". } \\
\\ 
\texttt{Next time, I will apply this law to simplify the expression before evaluating it. I will rewrite the expression as "not True" and then apply the NOT operator. This will give me the correct result without having to apply the NOT operator multiple times.} \\
\\ 
\texttt{Trial \#2: New plan: } \\
\\ 
\texttt{In my previous attempt, I correctly identified the double negation law and applied it to simplify the expression, but I failed to recognize that I could simplify the expression further by applying the law multiple times. However, I also failed to recognize that the expression "not not not not not not not True" can be simplified using the fact that "not True" is equivalent to "False". } \\
\\ 
\texttt{Next time, I will apply the double negation law to simplify the expression, and then I will recognize that "not False" is equivalent to "True". This will give me the correct result without having to apply the NOT operator multiple times.} \\
\\ 
\\ 
\texttt{Evaluate the result of a random Boolean expression.} \\
\\ 
\texttt{Q: not ( ( not not True ) ) is} \\
\texttt{A: False} \\
\\ 
\texttt{Q: True and False and not True and True is} \\
\texttt{A: False} \\
\\ 
\texttt{Q: not not ( not ( False ) ) is} \\
\texttt{A: True} \\
\\ 
\texttt{Q: not not not not not not not True is} \\
\texttt{A:} \\
            }
        }
    }
     \\ \\

    \theutterance \stepcounter{utterance}  
    & \multicolumn{4}{p{0.6\linewidth}}{
        \cellcolor[rgb]{0.561,0.737,0.561}{
            \makecell[{{p{\linewidth}}}]{
                \texttt{\tiny{[P1$\rangle$GM]}}
                \texttt{To evaluate the given Boolean expression "not not not not not not not True", I will follow the plan from Trial \#2.} \\
\\ 
\texttt{1. Apply the double negation law multiple times: } \\
\texttt{   "not not not not not not not True" can be simplified as "not not not not not True" } \\
\texttt{   which can be further simplified as "not not not not True" } \\
\texttt{   which can be further simplified as "not not not True" } \\
\texttt{   which can be further simplified as "not not True" } \\
\texttt{   which can be further simplified as "not True".} \\
\\ 
\texttt{2. Recognize that "not True" is equivalent to "False".} \\
\\ 
\texttt{Therefore, the result of the Boolean expression "not not not not not not not True" is False.} \\
            }
        }
    }
    & & \\ \\

    \theutterance \stepcounter{utterance}  
    & & & \multicolumn{2}{p{0.3\linewidth}}{
        \cellcolor[rgb]{0.9,0.9,0.9}{
            \makecell[{{p{\linewidth}}}]{
                \texttt{\tiny{[GM$|$GM]}}
                \texttt{Target: False} \\
            }
        }
    }
    & & \\ \\

    \theutterance \stepcounter{utterance}  
    & & & \multicolumn{2}{p{0.3\linewidth}}{
        \cellcolor[rgb]{0.9,0.9,0.9}{
            \makecell[{{p{\linewidth}}}]{
                \texttt{\tiny{[GM$|$GM]}}
                \texttt{game\_result = LOSE} \\
            }
        }
    }
    & & \\ \\

    \theutterance \stepcounter{utterance}  
    & & \multicolumn{4}{p{0.6\linewidth}}{
        \cellcolor[rgb]{0.95,0.95,0.95}{
            \makecell[{{p{\linewidth}}}]{
                \texttt{\tiny{[P1$\leftrightarrow$P]}}
                \texttt{Post{-}episode Reflexion (not sent to GM):} \\
\texttt{New plan: } \\
\\ 
\texttt{In my previous attempts, I correctly identified the double negation law and applied it to simplify the expression, but I failed to recognize that I could simplify the expression further by applying the law multiple times. However, I also failed to recognize that the expression "not not not not not not not True" can be simplified using the fact that "not True" is equivalent to "False". } \\
\\ 
\texttt{Next time, I will apply the double negation law to simplify the expression, and then I will recognize that "not False" is equivalent to "True". This will give me the correct result without having to apply the NOT operator multiple times.} \\
            }
        }
    }
    & \\ \\

\end{supertabular}
}

\end{document}
